\subsection{阿廷环的根}

命题 1. —— 设 A 是一个左阿廷环。A 的根 r 是 A 的最大幂零双边理想，且环 A/r 是半单的。

我们知道（VIII, 第 156 页，定理 1），A 的每个幂零双边理想都含于 r。我们来证明 r 是幂零的。由于环 A 是左阿廷的且双边理想的序列 \(r^{n}\) 是递降的，存在一个整数 \(p \geqslant 0\)，使得我们有 \(r^{p} = r^{p+1}\)。由于 A 是阿廷的，A-模 \(A_{s}\) 具有有限长度（VIII, 第 6 页，定理 1），因而是诺特的。于是左理想 \(r^{p}\) 是有限生成的。由中山引理（VIII, 第 158 页的定理 2）推出 \(r^{p} = 0\)。

环 A/r 无根（VIII, 第 155 页，命题 5）；由于它是左阿廷的，故它是半单的（VIII, 第 154 页，命题 4）。

推论 1. —— 一个环是半单的当且仅当它是左阿廷的且除 0 外没有幂零双边理想。

半单环是左阿廷的且无根（VIII, 第 154 页，命题 4）；因而它除 0 外没有幂零双边理想。逆命题由命题 1 得出。

推论 2. —— 在左阿廷环 A 中，根由满足对 A 中每个 a，ax 都是幂零的元素 x 组成。

若 A 是左阿廷的，则根是一个幂零双边理想（命题 1）。由雅各布森定理（VIII, 第 156 页，定理 1），每个左诣零理想都含于根中。推论 2 由此得出。

